\documentclass[11pt]{article}
\usepackage[localtheorems]{raeez-math-template}

\providecommand{\C}{\mathbb{C}}
\providecommand{\Q}{\mathbb{Q}}
\providecommand{\dbar}{\bar{\partial}}
\providecommand{\dbarstar}{\bar{\partial}^{\ast}}
\providecommand{\fg}{\mathfrak{g}}
\providecommand{\ad}{\operatorname{ad}}
\providecommand{\Tr}{\operatorname{Tr}}
\providecommand{\Sym}{\operatorname{Sym}}
\providecommand{\End}{\operatorname{End}}
\providecommand{\Rep}{\operatorname{Rep}}
\providecommand{\Alg}{\operatorname{Alg}}
\providecommand{\CoHA}{\operatorname{CoHA}}
\providecommand{\Obs}{\operatorname{Obs}}
\providecommand{\CE}{\operatorname{CE}}
\providecommand{\BV}{\operatorname{BV}}
\providecommand{\QME}{\operatorname{QME}}
\providecommand{\hCS}{\operatorname{hCS}}
\providecommand{\Ch}{\operatorname{Ch}}
\providecommand{\Dolb}{\operatorname{Dolb}}
\providecommand{\cA}{\mathcal{A}}
\providecommand{\cE}{\mathcal{E}}
\providecommand{\cF}{\mathcal{F}}
\providecommand{\cL}{\mathcal{L}}
\providecommand{\cO}{\mathcal{O}}
\providecommand{\cW}{\mathcal{W}}
\providecommand{\Eone}{E_1}
\providecommand{\Etwo}{E_2}
\providecommand{\Ethree}{E_3}
\providecommand{\PhiFA}{\Phi^{\mathrm{FA}}}
\providecommand{\SpCh}{\operatorname{SpCh}}
\providecommand{\Yp}{Y^+}
\providecommand{\Zchder}{Z^{\mathrm{der}}_{\mathrm{ch}}}
\providecommand{\kanom}{\kappa_{\mathrm{anom}}}
\providecommand{\kch}{\kappa_{\mathrm{ch}}}
\providecommand{\kcat}{\kappa_{\mathrm{cat}}}
\providecommand{\kBKM}{\kappa_{\mathrm{BKM}}}
\providecommand{\kfib}{\kappa_{\mathrm{fiber}}}

\theoremstyle{plain}
\newtheorem{theorem}{Theorem}[section]
\newtheorem{proposition}[theorem]{Proposition}
\newtheorem{lemma}[theorem]{Lemma}
\newtheorem{corollary}[theorem]{Corollary}
\theoremstyle{definition}
\newtheorem{definition}[theorem]{Definition}
\newtheorem{construction}[theorem]{Construction}
\newtheorem{hypothesis}[theorem]{Hypothesis}
\theoremstyle{remark}
\newtheorem{remark}[theorem]{Remark}

\title{Six-Dimensional Holomorphic Chern-Simons Observables:\\
The \(E_3\) Stage, BV Quantisation, and \(E_1\) Chiral Specialisation}
\author{}
\date{}

\begin{document}
\maketitle

\begin{abstract}
Holomorphic Chern-Simons theory on a complex Calabi-Yau threefold is a
Costello BV theory after fixing the holomorphic volume form, the
quadratic gauge coefficient, an elliptic gauge fixing, a heat-kernel
regulator, completed local functionals, and factorisation sewing. Its
local classical observables form the continuous Chevalley cochains of
the Dolbeault dg Lie algebra
\(\Omega^{0,\bullet}(-,\fg)[1]\). On a Stein polydisc in \(\C^3\), with
polynomial or compact support boundary conditions specified, the
holomorphic factorisation algebra is the \(E_3^{\mathrm{hol}}\)
enveloping object and satisfies the PBW filtration with associated
graded symmetric algebra on the shifted local Lie object. Quantum
statements on compact \(X\) require a solution of the scale-\(L\)
quantum master equation. In complex dimension three the one-loop
nonabelian gauge obstruction is the quartic class
\(I_4^\rho\in\Sym^4(\fg^\vee)^\fg\); the quadratic Casimir operator
\(C_2^\rho\) supplies the two-point field-strength counterterm. The
specialised \(d=3\) CY-to-chiral output on a reference curve is
\(E_1\)-chiral. Braided \(E_2\) structure appears on centre, double, or
module layers.
\end{abstract}

\tableofcontents

\section{Holomorphic Chern-Simons Data}
\label{sec:hcs-data}

A six-real-dimensional holomorphic Chern-Simons datum consists of
\[
  (X,\Omega_X,\fg,\langle-,-\rangle,h,Q^{\mathrm{GF}}),
\]
where \(X\) is a complex threefold with nowhere-vanishing holomorphic
volume form \(\Omega_X\), \(\fg\) is a finite-dimensional complex Lie
algebra with a non-degenerate invariant symmetric pairing, \(h\) is a
Hermitian metric, and \(Q^{\mathrm{GF}}\) is an elliptic Costello
gauge-fixing operator. The flat model is \(X=\C^3\) with
\(\Omega_X=dz_1\wedge dz_2\wedge dz_3\). Compact examples, such as
\(K3\times E\), need the additional quantisation package of
Section~\ref{sec:compact-package}.

\begin{definition}[BV fields]
\label{def:bv-fields}
For an open set \(U\subset X\), the shifted field complex is
\[
  \cE_U=\Omega^{0,\bullet}(U,\fg)[1].
\]
In the unshifted Dolbeault grading the \((0,0)\)-component is the
ghost, the \((0,1)\)-component is the physical field, and the
\((0,2)\)- and \((0,3)\)-components are BV antifields through the
pairing
\[
  \omega_{\BV}(\alpha,\beta)
  =
  \int_U\Omega_X\wedge\langle\alpha,\beta\rangle.
\]
The integral is taken on compactly supported variations, or on compact
\(U\). After the shift by \([1]\), this pairing has BV degree \(-1\).
\end{definition}

\begin{definition}[Classical action]
\label{def:classical-action}
The classical interaction on \(U\) is
\[
  S_{\mathrm{cl}}(a)=
  \int_U\Omega_X\wedge
  \left(
    \frac12\,\langle a,\dbar a\rangle
    +\frac16\,\langle a,[a,a]\rangle
  \right),
  \qquad
  a\in\Omega^{0,\bullet}_c(U,\fg)[1].
\]
The bracket is the wedge product of Dolbeault forms combined with the
Lie bracket on \(\fg\).
\end{definition}

\begin{proposition}[Classical master equation]
\label{prop:cme}
The functional \(S_{\mathrm{cl}}\) satisfies
\[
  \{S_{\mathrm{cl}},S_{\mathrm{cl}}\}_{\BV}=0.
\]
\end{proposition}

\begin{proof}
The Hamiltonian vector field of \(S_{\mathrm{cl}}\) is the Chevalley
differential of the dg Lie algebra
\[
  \bigl(\Omega^{0,\bullet}(U,\fg),\dbar,[-,-]\bigr).
\]
Its square is zero because \(\dbar^2=0\), \(\dbar\) is a derivation of
the wedge-Lie bracket, the Lie bracket satisfies Jacobi, and
\(\langle-,-\rangle\) is invariant. These identities are exactly the
three summands in the BV bracket of the kinetic and cubic terms.
\end{proof}

\begin{definition}[Local dg Lie algebra and classical observables]
\label{def:local-lie-observables}
Let
\[
  \cL_{\hCS}(U)=\Omega^{0,\bullet}_c(U,\fg)
\]
with differential \(\dbar\) and bracket
\([\alpha,\beta]=[\alpha\wedge\beta]_\fg\). The completed classical
observables are
\[
  \Obs^{\mathrm{cl}}_{\hCS}(U)
  =
  C^\bullet_{\CE,\mathrm{cont}}\bigl(\cL_{\hCS}(U)[1]\bigr)
  =
  \widehat{\Sym}\bigl(\cL_{\hCS}(U)[1]^\vee[-1]\bigr),
\]
with differential \(d_{\dbar}+d_{\CE}\).
\end{definition}

\section{The Local \(E_3\) Stage and PBW}
\label{sec:e3-pbw}

Holomorphic factorisation on \(\C^3\) is governed by configurations of
points in a complex threefold. It is an \(E_3^{\mathrm{hol}}\) structure
in \(\Ch(\Dolb)\), whose Dolbeault cohomology gives the ordinary
\(E_3\) operadic shadow in characteristic zero. It is a stage-one
factorisation structure. Section~\ref{sec:phi-specialisation} records
the drop to \(E_1\) after chiral specialisation.

\begin{theorem}[Local \(E_3\) PBW theorem]
\label{thm:e3-pbw}
Let \(U\subset\C^3\) be a Stein polydisc with compact-support
observables, or \(U=\C^3\) with polynomial boundary conditions. The
classical holomorphic Chern-Simons
factorisation algebra generated by \(\cL_{\hCS}(U)\) is the
\(E_3^{\mathrm{hol}}\)-enveloping factorisation algebra. Its
locally constant reduction has a PBW filtration \(F_\bullet\) with
\[
  \operatorname{gr}_F
  \mathrm{U}^{\mathrm{fact}}_{E_3}(\cL_{\hCS})
  \simeq
  \Sym_{\Ch(\Dolb)}\bigl(\cL_{\hCS}[2]\bigr)
\]
as an \(E_3\)-factorisation coalgebra. Equivalently,
\[
  \Obs^{\mathrm{cl}}_{\hCS}(U)
  =
  C^\bullet_{\CE,\mathrm{cont}}\bigl(\cL_{\hCS}(U)[1]\bigr)
\]
has associated graded the completed signed symmetric algebra on the
linear local fields.
\end{theorem}

\begin{proof}
Francis' \(E_n\) PBW theorem and Knudsen's higher enveloping theorem
identify the associated graded of the \(E_n\)-enveloping filtration
with the shifted symmetric algebra on the Lie input. For \(n=3\), the
shift is \([2]\), the degree of the Browder operation coming from
\[
  H^\bullet(\operatorname{Conf}_2(\mathbb{R}^3);\Q)
  \cong H^\bullet(S^2;\Q).
\]
The Dolbeault local Lie algebra \(\cL_{\hCS}\) is the input Lie object
for Costello-Gwilliam classical observables. Fulton-MacPherson
compactified configurations compute the symmetrisation map; boundary
strata encode iterated brackets, and Stokes' formula gives Jacobi with
the Koszul signs of Dolbeault forms.
\end{proof}

\begin{corollary}[Formal stalk]
\label{cor:formal-stalk}
At a formal point of \(\C^3\), the Dolbeault resolution contracts to
\(\C\). After completion and continuous duality,
\[
  \cL_{\hCS,x}\simeq \fg,
  \qquad
  \operatorname{gr}_F\mathrm{U}^{\mathrm{fact}}_{E_3}(\cL_{\hCS,x})
  \simeq
  \Sym(\fg[2]).
\]
A Milnor-Moore reconstruction of observables from \(E_3\)-primitives
requires augmentation, completion, conilpotence, and duality
hypotheses.
\end{corollary}

\begin{remark}[Boundary conditions]
\label{rem:boundary-conditions}
For polynomial boundary conditions on \(\C^3\),
\[
  H^\bullet_{\dbar}(\C^3,\fg)=\C[z_1,z_2,z_3]\otimes\fg
\]
in degree zero, with transferred bracket
\[
  [f\otimes x,g\otimes y]=fg\otimes[x,y]_\fg.
\]
For \(L^2\) boundary conditions on \(\C^3\), the harmonic subspace
vanishes. The boundary condition is part of the theorem.
\end{remark}

\begin{definition}[Supported observables]
\label{def:supported-observables}
The same stage-one factorisation object has different transverse
restrictions.
\begin{enumerate}
\item Point support in a three-complex-dimensional transverse slice
carries the \(E_3^{\mathrm{hol}}\) structure of
Theorem~\ref{thm:e3-pbw}.
\item A holomorphic curve \(C\subset X\) has a two-complex-dimensional
normal slice. The transverse restriction is controlled by the
\(E_2^{\mathrm{hol}}\) operations of that normal slice.
\item A holomorphic surface \(S\subset X\) has a one-complex-dimensional
normal slice. The transverse restriction is controlled by \(E_1\)
operations.
\end{enumerate}
These are support and normal-direction structures inside the same
factorisation algebra.
\end{definition}

\section{Heat-Kernel Propagator}
\label{sec:propagator}

\begin{definition}[Regularised propagator]
\label{def:propagator}
On \(\C^3\) with the Euclidean Hermitian metric and
\(Q^{\mathrm{GF}}=\dbarstar\), set
\[
  P_{\varepsilon,L}(z,w)
  =
  \int_\varepsilon^L(\dbarstar\otimes 1)K_t(z,w)\,dt,
  \qquad
  0<\varepsilon<L<\infty ,
\]
where
\[
  K_t(z,w)=(4\pi t)^{-3}\exp(-\|z-w\|^2/(4t))
\]
is the scalar heat kernel of the flat Dolbeault Laplacian, tensored
with the identity on the gauge coefficient.
\end{definition}

\begin{proposition}[Bochner-Martinelli representative]
\label{prop:bm}
\(P_{\varepsilon,L}\) is smooth on \(\C^3\times\C^3\). Off the
diagonal, the limit \(\varepsilon\to0\), \(L\to\infty\) is the
Bochner-Martinelli kernel
\[
  P_{\mathrm{BM}}(z,w)
  =
  \frac{2!}{(2\pi i)^3}
  \sum_{k=1}^{3}(-1)^{k-1}
  \frac{\overline{z_k-w_k}}{\|z-w\|^6}\,
  d\bar z_1\wedge\cdots\widehat{d\bar z_k}\cdots\wedge d\bar z_3
  \wedge dw_1\wedge dw_2\wedge dw_3.
\]
Its diagonal singularity has order \(\|z-w\|^{-5}\), and
\[
  \dbar_wP_{\mathrm{BM}}=\delta_\Delta
\]
as a distribution on \(\C^3\times\C^3\).
\end{proposition}

\begin{proof}
The flat Green operator for the Dolbeault Laplacian is
\(\int_0^\infty K_t\,dt\) off the diagonal. Applying
\(\dbarstar\) in one variable differentiates the Gaussian and produces
the numerator \(\overline{z_k-w_k}\). Integration over
\((\varepsilon,L)\) removes the ultraviolet diagonal singularity and
the infrared tail. Removing both cutoffs off the diagonal recovers the
standard Bochner-Martinelli representative. The distributional identity
is the Cauchy-Green identity for the \(\dbar\)-homotopy kernel.
\end{proof}

\begin{proposition}[Heat-kernel homotopy identity]
\label{prop:heat-homotopy}
The regularised kernel satisfies
\[
  \dbar P_{\varepsilon,L}=K_\varepsilon-K_L.
\]
Consequently a homotopy of Costello gauge fixings changes the
propagator by a \(Q\)-exact kernel. Once the QME obstruction vanishes,
Costello's independence theorem identifies the resulting quantum
observables up to equivalence of factorisation algebras.
\end{proposition}

\begin{proof}
The Dolbeault heat kernel satisfies
\(\partial_tK_t=-\Delta_{\dbar}K_t\), with
\(\Delta_{\dbar}=\dbar\dbarstar+\dbarstar\dbar\). Applying \(\dbar\) to
\((\dbarstar\otimes 1)K_t\) and integrating over
\([\varepsilon,L]\) telescopes to \(K_\varepsilon-K_L\), with the
remaining summand absorbed by the gauge-fixing homotopy.
\end{proof}

\section{The Quartic QME Obstruction}
\label{sec:qme-obstruction}

\begin{definition}[Quartic invariant]
\label{def:quartic}
For a representation \(\rho:\fg\to\End(V)\), define
\[
  I_4^\rho(x_1,x_2,x_3,x_4)
  =
  \frac1{4!}\sum_{\sigma\in S_4}
  \Tr_V\!\left(
    \rho(x_{\sigma(1)})\rho(x_{\sigma(2)})
    \rho(x_{\sigma(3)})\rho(x_{\sigma(4)})
  \right).
\]
Then \(I_4^\rho\in\Sym^4(\fg^\vee)^\fg\). The adjoint trace gives
\(I_4^{\ad}\).
\end{definition}

\begin{theorem}[One-loop QME obstruction in complex dimension three]
\label{thm:quartic-obstruction}
For nonabelian holomorphic Chern-Simons theory in complex dimension
three, the one-loop gauge contribution to the QME obstruction is
represented in the local Chevalley complex by
\[
  \Theta^\rho_U(\alpha_0,\alpha_1,\alpha_2,\alpha_3)
  =
  \int_U
  \left[
  I_4^\rho\bigl(
    \alpha_0,\partial\alpha_1,\partial\alpha_2,\partial\alpha_3
  \bigr)
  \right]_{(3,3)},
  \qquad
  \alpha_i\in\Omega^{0,\bullet}_c(U,\fg).
\]
Its class
\[
  \kanom(X,\fg,\rho)
  =
  [\,\hbar\,\Theta^\rho\,]
  \in
  H^1_{\mathrm{loc}}
  \bigl(\Omega^{0,\bullet}(X,\fg)[1],\cO_{\mathrm{loc}}\bigr)
\]
is the first obstruction to solving the scale-\(L\) QME for the
specified open hCS datum.
\end{theorem}

\begin{proof}
The one-loop wheel with four external gauge insertions supplies the
first degree-one local functional in complex dimension three. Its Lie
factor is the symmetrised trace of four representation matrices,
hence \(I_4^\rho\). Its holomorphic form factor contains three
\(\partial\)-derivatives; taking the \((3,3)\)-component produces a
local density on \(U\). Costello-Li identify the adjoint case with the
descent representative \(\int_X\Tr_{\ad}(cF_A^3)\), equivalently the
descent of \(\Tr_{\ad}(F_A^4)\). Costello's obstruction theory places
this local functional in degree one of the BV obstruction complex.
\end{proof}

\begin{definition}[Quadratic Casimir coefficient]
\label{def:casimir}
Let \(\{T_a\}\) be a basis of \(\fg\), and let \(\{T^a\}\) be the dual
basis for the chosen invariant pairing. The representation-theoretic
quadratic Casimir operator is
\[
  C_2^\rho=\sum_a\rho(T_a)\rho(T^a)\in\End_{\fg}(V).
\]
For the adjoint representation of a simple Lie algebra in the usual
normalisation, \(C_2^{\ad}=2h^\vee\,\mathrm{id}_{\fg}\).
\end{definition}

\begin{proposition}[Two-point field-strength counterterm]
\label{prop:field-strength}
With the flat heat-kernel normalisation, and when \(C_2^\rho\) acts by
the scalar \(c_2^\rho\) on the chosen representation block, the
two-point graph contributes a degree-zero local counterterm of the form
\[
  S_{\mathrm{ct}}^{(2)}
  =
  -\hbar\,\frac{c_2^\rho}{(4\pi)^3}
  \log(L/\varepsilon)
  \int_X\Omega_X\wedge\Tr_\rho(\cA\,\dbar\cA),
\]
up to the convention for the invariant pairing and the heat-kernel
normalisation. This term renormalises the kinetic pairing and the
field \(\cA\). The degree-one QME obstruction is the quartic class of
Theorem~\ref{thm:quartic-obstruction}.
\end{proposition}

\begin{proof}
Closing two legs in the cubic interaction yields the colour contraction
\(\sum_a\rho(T_a)\rho(T^a)=C_2^\rho\). On an irreducible or fixed
central representation block this operator acts by \(c_2^\rho\). The
remaining external fields have the tensor type of the kinetic functional
\(\int_X\Omega_X\wedge\Tr_\rho(\cA\dbar\cA)\). The short-distance
six-real-dimensional heat-kernel integral supplies the logarithmic
factor. Since the resulting local functional has cohomological degree
zero, it changes the kinetic representative. The arity-four wheel of
Theorem~\ref{thm:quartic-obstruction} occupies the obstruction slot.
\end{proof}

\begin{corollary}[Flat and abelian cases]
\label{cor:flat-abelian}
If \(\fg\) is abelian, the cubic interaction vanishes and the theory is
Gaussian. If the chosen representation satisfies \(I_4^\rho=0\), or if
\(\Theta^\rho\) is exact after the specified counterterm or closed-sector
package is included, the one-loop gauge obstruction vanishes for that
datum. Higher quantum statements are obtained by the same Costello
induction after the corresponding higher local obstruction classes
vanish or are trivialised.
\end{corollary}

\section{Compact Calabi-Yau Threefolds}
\label{sec:compact-package}

\begin{hypothesis}[Compact BV package]
\label{hyp:compact-bv}
A compact nonabelian hCS quantisation theorem on a Calabi-Yau
threefold requires the following data.
\begin{enumerate}
\item A compact Calabi-Yau threefold \((X,\Omega_X)\), together with an
orientation and determinant-line framing compatible with the BV
pairing.
\item A quadratic gauge Lie algebra or dg Lie algebra
\((\fg,\langle-,-\rangle)\), and a representation \(\rho\) whenever
trace invariants are used.
\item An elliptic Costello gauge fixing \(Q^{\mathrm{GF}}\), heat-kernel
regulator \(P_{\varepsilon,L}\), and RG-compatible counterterm scheme.
\item A completed space of local functionals and observables on which
\(\Delta_L\), \(\{-,-\}_L\), and RG flow are defined.
\item Factorisation or TCFT sewing data with descent from local
polydisks to \(X\).
\item A trivialisation of \(\kanom(X,\fg,\rho)\), by vanishing in local
BV cohomology, an explicit local counterterm, or an open-closed
Green-Schwarz term.
\end{enumerate}
\end{hypothesis}

\begin{theorem}[Conditional compact quantisation]
\label{thm:conditional-compact}
Assume Hypothesis~\ref{hyp:compact-bv}. If
\(\kanom(X,\fg,\rho)=0\) in the Costello local BV obstruction complex
after the chosen open or open-closed counterterms are included, then
the effective interactions \(I[L]\) can be chosen to satisfy the QME
through one loop. Under the analogous higher-loop vanishing or
counterterm hypotheses, the completed observables form a quantum
factorisation algebra on \(X\). A surviving obstruction class is
exactly the failure of the QME for the specified regulator and field
content.
\end{theorem}

\begin{proof}
Costello's obstruction theory solves the QME inductively in powers of
\(\hbar\). At the first quantum step, the obstruction is the
degree-one local class of Theorem~\ref{thm:quartic-obstruction},
modified by any local or closed-sector counterterm contribution.
Vanishing gives the first quantum correction. The same induction
applies at each later loop order inside the completed local functional
complex fixed in Hypothesis~\ref{hyp:compact-bv}.
\end{proof}

\begin{remark}[Characteristic numbers and anomaly data]
\label{rem:characteristic-numbers}
The hCS anomaly invariant \(\kanom(X,\fg,\rho)\) is separate from the
Calabi-Yau numerical invariants used elsewhere in the manuscript:
\[
  \kch,\qquad
  \kcat=\chi(\cO_X),\qquad
  \kBKM=c_N(0)/2,\qquad
  \kfib.
\]
For \(X=K3\times E\),
\[
  \chi_{\mathrm{top}}(K3\times E)=24\cdot0=0,
  \qquad
  \kcat(K3\times E)=\chi(\cO_{K3})\chi(\cO_E)=2\cdot0=0.
\]
The standard Vol III comparison values are
\[
  \kch^{\mathrm{compact}}(K3\times E)=0,\qquad
  \kch^{\mathrm{Heis}}(K3\times E)=3,\qquad
  \kBKM(\mathfrak g_{\Delta_5})=5,\qquad
  \kfib(K3)=24.
\]
They come from distinct constructions: compact Hodge supertrace,
Heisenberg-Mukai specialisation, Gritsenko's \(\Delta_5\) Borcherds
denominator, and the K3 Mukai lattice rank. The quartic hCS anomaly
calculation supplies a gauge-theoretic obstruction class. Any relation
between \(\kanom\) and the four comparison constants requires a
separate CY-to-chiral comparison theorem with compact QME and
Hall-Borcherds data.
\end{remark}

\section{Specialisation to a Chiral Algebra}
\label{sec:phi-specialisation}

The stage-one \(E_3^{\mathrm{hol}}\) structure is the local
factorisation object. The \(d=3\) CY-to-chiral output is obtained by
specialising along a complex surface and then restricting to a
reference curve.

\begin{proposition}[\(d=3\) output level]
\label{prop:d3-output}
Let \(\cC\) be a CY\(_3\) category on a witnessed hCS or Hochschild
comparison locus, and let \((\Sigma_2,C)\) be a witnessed admissible
specialisation datum. The two-stage construction is
\[
  \Phi_3^{(\Sigma_2,C)}(\cC)
  =
  \SpCh_{\Sigma_2,C}\bigl(\PhiFA_3(\cC)\bigr)
  =
  \left(\int_{\Sigma_2}^{\mathrm{hol}}\PhiFA_3(\cC)\right)\Big|_C
  \in
  \Alg^{\mathrm{ch}}_{\Eone}(C).
\]
The \(E_3^{\mathrm{hol}}\) structure belongs to the stage-one
factorisation algebra. The braided \(E_2\) structure belongs to centre,
double, or module layers. On loci where the derived chiral centre and
the representation-centre comparison are constructed, one has
\[
  Z\bigl(\Rep^{\Eone}(A_C)\bigr)
  \simeq
  \Rep^{\Etwo}\bigl(\Zchder(A_C)\bigr)
\]
on the verified Drinfeld-centre locus.
\end{proposition}

\begin{proof}
The stage-one object \(\PhiFA_3(\cC)\) is an \(E_3\)-holomorphic
factorisation algebra after the formality, framing, anomaly,
completion, and comparison witnesses have been fixed. Factorisation
homology over the complex surface \(\Sigma_2\) lowers the complex
operadic level by two. Restriction to \(C\) therefore leaves an
\(E_1\)-chiral algebra on the curve. The Ben-Zvi-Francis-Nadler
centre mechanism recovers \(E_2\)-braiding on the representation
centre when its hypotheses are satisfied.
\end{proof}

\begin{remark}[\(\C^3\) CoHA normalisation]
\label{rem:c3-coha}
For the toric Hall algebra of \(\C^3\),
\[
  \CoHA(\C^3)\simeq \Yp(\widehat{\mathfrak{gl}}_1),
\]
the positive half in the Schiffmann-Vasserot normalisation. The full
affine Yangian is the Drinfeld double
\[
  D\bigl(\Yp(\widehat{\mathfrak{gl}}_1)\bigr)
  =
  \widehat Y(\widehat{\mathfrak{gl}}_1).
\]
The vertex algebra \(\cW_{1+\infty}\) appears after the corresponding
double or centre, completion, and evaluation representation are
constructed. PBW for hCS observables and PBW for
\(\Yp(\widehat{\mathfrak{gl}}_1)\) agree only on the constructed toric
comparison locus.
\end{remark}

\section{Separated Objects}
\label{sec:separated-objects}

\begin{remark}[Fields, antifields, and categorical deformation data]
\label{rem:fields-vs-polyvectors}
The hCS fields are \(\Omega^{0,\bullet}(X,\operatorname{ad}P)[1]\).
Their BV antifields are the complementary Dolbeault components paired
by \(\int_X\Omega_X\wedge\Tr_\rho(-\, -)\). Hochschild polyvectors of
a CY category are a different object:
\[
  \mathrm{HH}^\bullet(\operatorname{Perf}(\C^3))
  \simeq
  \C[x,y,z]\otimes
  \bigwedge\nolimits^\bullet
  \langle\partial_x,\partial_y,\partial_z\rangle.
\]
These polyvectors control intrinsic categorical deformations and the
derived centre. A finite-dimensional hCS gauge algebra enters after a
framing object or external gauge coefficient has been chosen.
\end{remark}

\begin{remark}[Five algebraic objects]
\label{rem:five-objects}
For an \(E_1\)-chiral algebra \(A\), the objects
\[
  A,\qquad B(A),\qquad A^i,\qquad A^!,\qquad \Zchder(A)
\]
have different definitions and functorial roles. The equivalence
\(\Omega B(A)\simeq A\) is bar-cobar inversion. The object \(A^i\) is
the Koszul dual coalgebra. The object \(A^!\) is the Verdier or
Koszul dual algebra, depending on the duality context. The bulk or
braided layer at \(d\geq3\) is \(\Zchder(A)\), the derived chiral
centre on the constructed loci.
\end{remark}

\section{References}
\label{sec:references}

The BV renormalisation formalism is Costello, \emph{Renormalization
and Effective Field Theory}, 2011. The factorisation-algebra
construction of classical and quantum observables is Costello and
Gwilliam, \emph{Factorization Algebras in Quantum Field Theory},
Volumes 1 and 2, 2017 and 2021. The holomorphic anomaly and
open-closed cancellation mechanism are Costello and Li's holomorphic
BV analysis of the B-model. The \(E_n\) PBW input is Francis 2013 and
Knudsen 2018. Holomorphic factorisation algebras and their
\(E_n^{\mathrm{hol}}\) description are treated by Gwilliam and
Williams. The toric CoHA normalisation
\(\CoHA(\C^3)\simeq\Yp(\widehat{\mathfrak{gl}}_1)\) is the
Kontsevich-Soibelman and Schiffmann-Vasserot positive-half
normalisation; \(\cW_{1+\infty}\) enters through the subsequent
double, centre, completion, and evaluation chain.

\end{document}
